\documentclass[11pt]{article}

\usepackage[margin=1in]{geometry}
\usepackage[T1]{fontenc}
\usepackage[utf8]{inputenc}
\usepackage{amsmath,amssymb}
\usepackage{booktabs}
\usepackage{enumitem}
\usepackage{graphicx}
\usepackage{hyperref}
\usepackage{longtable}
\usepackage{tabularx}
\usepackage{xcolor}
\usepackage{xurl}

\hypersetup{
    colorlinks=true,
    linkcolor=blue!55!black,
    urlcolor=blue!55!black,
    citecolor=blue!55!black
}

\setlist[itemize]{leftmargin=*}
\setlist[enumerate]{leftmargin=*}
\renewcommand{\arraystretch}{1.15}
\graphicspath{{../}{./}}

\newcommand{\deePC}{DeePC}
\newcommand{\prr}{\mathrm{PRR}}
\newcommand{\pir}{\mathrm{PIR}}
\newcommand{\cbr}{\mathrm{CBR}}
\newcommand{\artifact}[1]{\path{#1}}

\title{\textbf{Progress Evaluation Report: \deePC-Based NR-V2X Sidelink Transmit-Power Adaptation}}
\author{Bim\\
\small ns-3 NR-V2X / NGSIM US-101 DeePC Study}
\date{July 2026}

\begin{document}

\maketitle

\begin{center}
\begin{tabular}{ll}
\toprule
Report purpose & Professor evaluation of completed research and engineering progress \\
Main deliverable & Written progress-evaluation report \\
Claim status & Feasibility and implementation progress; final performance claims pending \\
\bottomrule
\end{tabular}
\end{center}

\tableofcontents
\newpage

\section{Executive Summary}

This report summarizes the current state of my research on Data-Enabled
Predictive Control (\deePC) for adaptive NR-V2X sidelink communication. The
central goal is to test whether measured vehicle-to-everything input-output
data can support online transmit-power control in a realistic ns-3 NR-V2X
simulation, with the controller balancing communication reliability,
timeliness, and congestion. Reliability is measured using awareness-range
packet reception ratio (PRR), timeliness is measured using packet
inter-reception time (PIR), and congestion is measured using channel busy ratio
(CBR).

The main completed milestone is an end-to-end \deePC{} research pipeline. I
have constructed \deePC-ready datasets from ns-3 KPI traces, generated
normalized Hankel matrices, performed offline data-window checks, built
closed-loop controller infrastructure through an ns-3 JSON file bridge, added a
Python OSQP controller for online transmit-power control, prepared
MATLAB/YALMIP export support, and generated campaign-level analysis tables and
readiness reports. The current controller can receive recent ns-3 measurements,
solve the \deePC{} optimization problem, and return a transmit-power command
during the simulation.

The strongest current result is feasibility rather than final performance
superiority. In the available preliminary closed-loop run, the \deePC{}
controller completed 285 online control updates with a 100\% controller
success rate and a mean solve time of about 0.0199 s, which is far below the
0.5 s control interval. The preliminary KPI values are close to the fixed
10 Hz baseline: mean PRR is 0.5739 for \deePC{} and 0.5748 for fixed 10 Hz in
the currently available partial evidence. Therefore, the work currently
supports the claim that the closed-loop \deePC{} control pipeline is
implemented and real-time compatible in this simulation setup. It does not yet
support a final claim that \deePC{} improves PRR over fixed 10 Hz operation.

The remaining major evaluation step is to complete matched repeated ns-3 runs
for the main methods. The current paper-readiness report shows that the
intended repeated campaign is not complete: fixed 10 Hz has partial run-101
evidence, the \deePC{} run used for orientation is an extra non-matched run,
and run IDs 101--105 are still missing for the main \deePC{} method. For this
reason, all closed-loop comparison results in this report are treated as
preliminary feasibility evidence.

\subsection*{Evaluation Focus}

The main points for evaluation are whether the ns-3 NR-V2X data pipeline is
valid, whether the \deePC{} formulation is correctly constructed from measured
input-output data, whether the ns-3 file bridge enables closed-loop control,
whether the Python OSQP controller can solve within the 0.5 s control interval,
and whether the preliminary 10 Hz results justify completing the matched
repeated-run campaign.

\section{Problem and Motivation}

NR-V2X sidelink communication enables vehicles to exchange cooperative
awareness messages directly, without routing every message through cellular
infrastructure. This is important for cooperative and automated driving because
nearby vehicles need timely information about position, speed, and motion
state. In dense traffic, however, periodic message generation can create
resource contention and interference. If transmit power is too low, vehicles
may miss important packets. If transmit power is unnecessarily high, the
channel can become more congested and neighboring transmissions may interfere
more strongly.

Fixed communication settings are easy to reproduce and are useful as baselines,
but they do not react to measured network conditions. A fixed 10 Hz CAM
configuration may be reasonable in moderate conditions, yet the best
transmit-power choice can depend on the number of active vehicles, recent PRR,
recent CBR, and other measured state. The research question for this stage is:

\begin{quote}
\textbf{Can measured NR-V2X input-output data support real-time predictive
transmit-power control using \deePC?}
\end{quote}

This question is deliberately narrower than global optimization of NR-V2X. The
current controller does not change SPS behavior, MCS, sensing parameters, CAM
payload structure, or the reservation period. It focuses on transmit power as
the controlled input while keeping the main comparison at fixed 10 Hz CAM
generation. This makes the comparison with a fixed 10 Hz baseline cleaner:
both methods generate CAMs at the same nominal rate, while \deePC{} is allowed
to adapt transmit power.

The control objective is to maintain or improve PRR while avoiding excessive
CBR and keeping PIR reasonable. PRR and PIR measure application-facing
awareness quality, while CBR measures congestion at the PHY layer. These
metrics create a natural tradeoff: a lower message rate or lower transmit
power can reduce CBR, but it may also increase PIR or reduce awareness
reliability. \deePC{} is attractive because it can use collected input-output
data directly rather than requiring a closed-form analytical model of the full
NR-V2X network.

\section{System and Dataset}

The simulation environment is ns-3 with an NR-V2X sidelink scenario driven by
processed NGSIM US-101 highway mobility. The final campaign folder identifies
the scenario as ``NR-V2X CAM \deePC{} control over NGSIM US-101 mainline
active-window mobility'' and uses the mobility file
\artifact{data/processed/ngsim_us101_mainline_active20_250_densest_600s.csv}.
The campaign configuration uses a 250-vehicle physical UE pool and dynamic
active assignment. The processed mobility metadata reports 852 logical vehicle
tracks, a peak of 179 simultaneous active vehicles, and a peak of 118 active
core vehicles.

\begin{table}[htbp]
\centering
\caption{Main evaluation configuration.}
\label{tab:configuration}
\begin{tabularx}{\linewidth}{lX}
\toprule
Item & Value \\
\midrule
Simulator & ns-3 NR-V2X sidelink \\
Mobility & NGSIM US-101 processed highway trajectories \\
Physical UE pool & 250 UEs \\
Simulation target & 300 s \\
Warmup / cooldown & 10 s / 10 s \\
KPI sample time & 0.5 s for final campaign data \\
Control interval & 0.5 s \\
Awareness range & 200 m \\
Main CAM interval & 0.1 s, or 10 Hz \\
Carrier frequency & 5.9 GHz \\
MCS & 6 \\
HARQ & Disabled in the current campaign \\
Sidelink sensing & Enabled \\
\deePC{} transmit-power bounds & 10 to 23 dBm \\
\bottomrule
\end{tabularx}
\end{table}

\begin{table}[htbp]
\centering
\caption{Main \deePC{} signals.}
\label{tab:signals}
\begin{tabularx}{\linewidth}{lX}
\toprule
Role & Signals \\
\midrule
Controlled input & \texttt{tx\_power\_dbm} \\
Fixed/context input & \texttt{beacon\_interval\_s} \\
Outputs & \texttt{prr\_awareness}, \texttt{pir\_s}, \texttt{cbr} \\
Traffic context & \texttt{active\_vehicle\_count\_core} and related density variables \\
\bottomrule
\end{tabularx}
\end{table}

The primary closed-loop dataset for the final campaign is stored under
\artifact{data/output/final/ieee_250veh_deepc_campaign_01/01_training/}. The
Tx-power-only, fixed-10-Hz \deePC{} dataset is
\artifact{deepc_dataset_txpower_count_tb0p1_dt0p5_n10}. Its source KPI file is
\artifact{kpi_timeseries_tb0p1_for_deepc.csv}, and it uses a sample time of
0.5 s. The clean evaluation window contains 561 rows from 10.0 s to 290.0 s.
The split is 336 training rows, 112 validation rows, and 113 test rows. The
\deePC{} past horizon is 20 samples, or 10 s, and the future horizon is
10 samples, or 5 s.

\begin{table}[htbp]
\centering
\caption{Hankel blocks for the Tx-power-only final-campaign controller.}
\label{tab:hankel}
\begin{tabular}{lc}
\toprule
Hankel block & Shape \\
\midrule
\texttt{u\_p} & $20 \times 307$ \\
\texttt{u\_f} & $10 \times 307$ \\
\texttt{y\_p} & $60 \times 307$ \\
\texttt{y\_f} & $30 \times 307$ \\
\texttt{d\_p} & $40 \times 307$ \\
\texttt{d\_f} & $20 \times 307$ \\
\bottomrule
\end{tabular}
\end{table}

\subsection{Collected-Data Validity for Hankel Construction}

Before constructing the Hankel matrices, I checked the collected ns-3 KPI
trace for internal consistency and input excitation coverage. These checks do
not represent field-measurement validation, but they are important for
confirming that the \deePC{} matrices are built from coherent simulator-side
input-output data.

Figure~\ref{fig:packet-accounting} checks the packet-accounting basis of the
training data. The PRR-related outputs are derived from transmitted packets,
awareness-eligible packet opportunities, and successful receptions over time.
This supports the use of the collected KPI trace as the measured output
sequence for \deePC{} training.

\begin{figure}[htbp]
    \centering
    \includegraphics[width=0.92\linewidth]{hankel_data_quality_corrected/packet_accounting_timeseries.png}
    \caption{Packet-accounting time series used to validate the collected KPI traces before constructing Hankel matrices.}
    \label{fig:packet-accounting}
\end{figure}

Figure~\ref{fig:prbs-inputs} shows the input trajectory and traffic context
used before Hankel matrix construction. The transmit-power input changes over
the training horizon, while the active-vehicle count captures time-varying
traffic density. This variation is important because \deePC{} relies on
informative input-output trajectories rather than a hand-derived analytical
network model.

\begin{figure}[htbp]
    \centering
    \includegraphics[width=0.92\linewidth]{hankel_data_quality_corrected/prbs_inputs_timeseries.png}
    \caption{Input excitation and traffic-context time series used before Hankel matrix construction.}
    \label{fig:prbs-inputs}
\end{figure}

\section{DeePC Mathematical Formulation}

This section describes the mathematical controller used in the closed-loop
Python OSQP implementation. The formulation is data-driven: it does not use an
analytical NR-V2X channel or queueing model. Instead, it uses previously
collected input-output trajectories from ns-3 to construct Hankel matrices and
then chooses the future transmit-power sequence that best matches the desired
communication behavior.

\subsection{Signals and Normalization}

At each sample $k$, the controller uses one control input,
\[
u_k =
\begin{bmatrix}
\texttt{tx\_power\_dbm}
\end{bmatrix},
\]
three measured outputs,
\[
y_k =
\begin{bmatrix}
\texttt{prr\_awareness} &
\texttt{pir\_s} &
\texttt{cbr}
\end{bmatrix}^{\top},
\]
and context variables,
\[
d_k =
\begin{bmatrix}
\texttt{active\_vehicle\_count\_core} &
\texttt{beacon\_interval\_s}
\end{bmatrix}^{\top}.
\]
For the main 10 Hz controller, the beacon interval is treated as fixed context
at 0.1 s, while transmit power is the optimized input. All input, output, and
context signals are normalized using the means and standard deviations saved in
\artifact{scaler.json}. The optimization is solved in normalized coordinates,
and the selected transmit power is converted back to dBm before being returned
to ns-3.

\subsection{Hankel Data Representation}

Let the past horizon be $T_{\mathrm{ini}} = 20$ samples and the future horizon
be $N = 10$ samples. Since the KPI/control sample time is 0.5 s, these
correspond to 10 s of history and a 5 s prediction horizon. From the normalized
training trajectory, block Hankel matrices are built and partitioned into past
and future blocks:
\[
\begin{bmatrix}
U_p \\
U_f \\
Y_p \\
Y_f \\
D_p \\
D_f
\end{bmatrix}.
\]
For the Tx-power-only final-campaign controller, the exported dimensions are:
\[
U_p \in \mathbb{R}^{20 \times 307}, \quad
U_f \in \mathbb{R}^{10 \times 307},
\]
\[
Y_p \in \mathbb{R}^{60 \times 307}, \quad
Y_f \in \mathbb{R}^{30 \times 307},
\]
\[
D_p \in \mathbb{R}^{40 \times 307}, \quad
D_f \in \mathbb{R}^{20 \times 307}.
\]
The 307 columns correspond to the number of available Hankel windows from the
training split. The decision variable is the coefficient vector
\[
g \in \mathbb{R}^{307},
\]
which selects a linear combination of previously observed trajectory segments.
The central DeePC assumption is that a future trajectory consistent with the
measured system can be represented through these Hankel columns.

\subsection{Past Matching and Context Forecast}

During closed-loop simulation, ns-3 sends the controller the most recent
history vectors:
\[
u_{\mathrm{ini}} \in \mathbb{R}^{20}, \quad
y_{\mathrm{ini}} \in \mathbb{R}^{60}, \quad
d_{\mathrm{ini}} \in \mathbb{R}^{40}.
\]
The future context vector $d_f^{\mathrm{ref}}$ is not optimized. It is supplied
as a forecast: the current implementation holds the latest active-vehicle count
constant over the 10-step horizon and fixes the beacon interval at 0.1 s. The
trajectory coefficient $g$ must match the measured past and the forecast
context:
\[
U_p g = u_{\mathrm{ini}},
\]
\[
Y_p g = y_{\mathrm{ini}},
\]
\[
D_p g = d_{\mathrm{ini}}, \qquad D_f g = d_f^{\mathrm{ref}}.
\]
The predicted future input and output trajectories are then
\[
\hat{u}_f = U_f g,
\]
\[
\hat{y}_f = Y_f g.
\]
Because $y_k$ has three components and $N=10$, $\hat{y}_f$ contains 30
normalized values: predicted PRR, PIR, and CBR for each future step.

\subsection{Optimization Problem}

The Python controller solves the following quadratic program at each control
update:
\[
\begin{aligned}
\min_{g} \quad
& \sum_{i=1}^{N}
    \left\lVert \hat{y}_{i} - y^{\mathrm{ref}} \right\rVert_Q^2
  + \sum_{i=1}^{N}
    \left\lVert \hat{u}_{i} \right\rVert_R^2
  + \lambda_g \lVert g \rVert_2^2 \\
\text{s.t.} \quad
& U_p g = u_{\mathrm{ini}}, \\
& Y_p g = y_{\mathrm{ini}}, \\
& D_p g = d_{\mathrm{ini}}, \\
& D_f g = d_f^{\mathrm{ref}}, \\
& u_{\min} \leq U_f g \leq u_{\max}.
\end{aligned}
\]
Here, $y^{\mathrm{ref}}$ is the desired output target. The current controller
uses:
\[
y^{\mathrm{ref}} =
\begin{bmatrix}
0.90 & 0.10 & 0.35
\end{bmatrix}^{\top},
\]
corresponding to high PRR, low PIR, and moderate CBR. The output tracking
weights are:
\[
Q = \mathrm{diag}(40,\;3,\;8),
\]
so PRR tracking is weighted most strongly, followed by CBR and PIR. The input
penalty is $R = 0.002$, and the coefficient regularization is
$\lambda_g = 10$. The physical transmit-power bounds are:
\[
10 \leq \texttt{tx\_power\_dbm} \leq 23.
\]
In the implementation these bounds are normalized before being passed to OSQP.

\subsection{Receding-Horizon Application}

After OSQP solves the quadratic program, the controller obtains the optimized
future transmit-power sequence
\[
\hat{u}_f =
\begin{bmatrix}
\hat{u}_{k|k} &
\hat{u}_{k+1|k} &
\cdots &
\hat{u}_{k+N-1|k}
\end{bmatrix}^{\top}.
\]
Only the first value $\hat{u}_{k|k}$ is sent back to ns-3 and applied during
the next 0.5 s control interval. At the next update, ns-3 provides a new
measured history, the optimization is solved again, and the process repeats.
This is the standard receding-horizon structure. The controller also logs the
predicted PRR, PIR, CBR horizon, solve time, solver status, and equality
residuals for later validation.

\subsection{Interpretation for NR-V2X}

In this formulation, \deePC{} is not predicting packets directly. It predicts
future KPI behavior using measured input-output trajectories. The optimization
therefore chooses transmit power by balancing three goals: keep PRR high, keep
PIR low, and avoid excessive CBR. The current implementation controls only
transmit power, while the 10 Hz CAM generation interval remains fixed. This is
why the main closed-loop comparison is against the fixed 10 Hz baseline.

\section{Implemented Work}

\subsection{DeePC Data Pipeline}

The first completed milestone was the offline \deePC{} data pipeline. I
created a configurable script that can clean the KPI series, select
input/output/context columns, split data into train/validation/test segments,
normalize each signal, and export Hankel matrices. The generated artifacts
include clean CSVs, normalized CSVs, \artifact{scaler.json},
\artifact{dataset_summary.json}, and \artifact{hankel_train_normalized.npz}.

\subsection{Offline Data-Window Check}

The open-loop ns-3 run itself does not use \deePC{} control. It is only the
data-collection run used to record input, output, and context trajectories.
After that run, I constructed the \deePC{} dataset and Hankel matrices from the
recorded data. As a separate offline diagnostic, I also fitted a linear ridge
map from past input/output/context windows and known future input/context
windows to future PRR, PIR, and CBR windows. This diagnostic checks whether the
recorded open-loop trace contains usable input-output structure before the data
is used in the closed-loop \deePC{} controller. The diagnostic metrics are
shown in Table~\ref{tab:open-loop}.

\begin{table}[htbp]
\centering
\caption{Offline held-out data-window check from the recorded open-loop ns-3 dataset.}
\label{tab:open-loop}
\begin{tabular}{llrrr}
\toprule
Diagnostic & Output & RMSE & MAE & $R^2$ \\
\midrule
Linear ridge window model & PRR & 0.049929 & 0.038993 & -1.063849 \\
Linear ridge window model & PIR & 0.039214 & 0.032651 & 0.113547 \\
Linear ridge window model & CBR & 0.019695 & 0.016011 & 0.356386 \\
\bottomrule
\end{tabular}
\end{table}

\noindent\footnotesize Source:
\artifact{data/output/final/ieee_250veh_deepc_campaign_01/01_training/open_loop_predictions/metrics.json}.
\normalsize

These RMSE values are not closed-loop performance results, and \deePC{} was
not active during the open-loop ns-3 simulation. They are a post-processing
diagnostic used before closed-loop control to check whether the collected
input-output trace is suitable for Hankel-matrix construction. The current
diagnostic is strongest for CBR and weaker for PRR, so PRR-related performance
claims should come from the matched closed-loop campaign rather than from this
offline check.

\subsection{Closed-Loop Infrastructure}

The ns-3 scenario has been prepared for true controller-in-the-loop simulation
through a JSON file bridge. In bridge mode, ns-3 writes request files
containing recent control, KPI, and context histories. An external controller
reads those requests, solves an optimization problem, writes response files,
and ns-3 applies the returned control action.

Implemented bridge behavior includes:
\begin{itemize}
    \item request files under \artifact{bridge/requests/request_XXXXXX.json};
    \item response files under \artifact{bridge/responses/response_XXXXXX.json};
    \item applied-control logging in \artifact{bridge/applied_controls.csv};
    \item controller timing in \artifact{bridge/controller_solve_times.csv};
    \item protection against overwriting the original open-loop dataset;
    \item rejection of bridge mode when the PRBS input schedule is still enabled.
\end{itemize}

\subsection{Controllers}

Several controller paths have been implemented or prepared. First, a
MATLAB/YALMIP path was prepared by exporting the normalized Hankel matrices and
metadata to \artifact{matlab_deepc_data.mat}. The MATLAB file controller can
watch the same bridge directory, solve a classical \deePC{} optimization in
YALMIP, and write response JSON files using the same schema. This path has not
been executed inside the Linux workspace because MATLAB is available separately
on Windows.

Second, a Python OSQP file controller was implemented in
\artifact{scripts/deepc_osqp_file_controller.py}. This controller loads
\artifact{dataset_summary.json}, \artifact{scaler.json}, and
\artifact{hankel_train_normalized.npz}, checks that the request horizons and
signal columns match the dataset, builds a CVXPY/OSQP optimization problem,
enforces transmit-power bounds from 10 to 23 dBm, and writes predicted PRR,
PIR, CBR, future transmit-power trajectories, solver status, solve time, and
equality-residual diagnostics.

Third, an acados-based offline surrogate controller was implemented. The
acados work verifies that a receding-horizon controller can be generated and
rolled out on a learned surrogate. The surrogate rollout improved predicted PRR
and lowered predicted CBR relative to the measured held-out open-loop segment,
but this is not yet an ns-3 closed-loop result and should not be presented as
final network performance evidence.

\subsection{Campaign Tooling}

The final campaign folder is
\artifact{data/output/final/ieee_250veh_deepc_campaign_01/}. It contains a
campaign manifest, training artifacts, repeated-run directories, command
scripts, analysis tables, plots, and an acceptance checklist. The aggregation
script \artifact{scripts/aggregate_final_campaign.py} produces paper-facing
outputs including \artifact{paper_readiness_report.md},
\artifact{paper_main_table.csv}, \artifact{aggregate_metrics.csv},
\artifact{run_metrics.csv}, density-bin tables, and comparison plots.

This is important because the project is now reproducible enough to evaluate:
there is a documented folder contract, exact command scripts for major stages,
and generated readiness reports that identify which claims are safe and which
runs are still missing.

\section{Code and Artifact Map}

To make the implementation easier to inspect, Table~\ref{tab:code-map} maps
the main reported claims to the script or artifact that supports them. The
Python OSQP controller source has also been commented around the bridge
protocol, signal validation, normalization, Hankel-block usage, optimization
objective, constraints, diagnostics, and receding-horizon response path.

\begin{longtable}{p{0.34\linewidth}p{0.59\linewidth}}
\caption{Main implementation files and artifacts for professor review.}
\label{tab:code-map}\\
\toprule
File or artifact & Role in the project \\
\midrule
\endfirsthead
\toprule
File or artifact & Role in the project \\
\midrule
\endhead
\artifact{scripts/build_hankel_dataset.py} & Builds the cleaned \deePC{} dataset, train/validation/test splits, normalization metadata, and Hankel matrices. \\
\artifact{scripts/open_loop_deepc_predict.py} & Runs the offline held-out data-window diagnostic reported in Table~\ref{tab:open-loop}. \\
\artifact{scripts/deepc_osqp_file_controller.py} & Implements the Python CVXPY/OSQP \deePC{} controller used by the ns-3 file bridge in the closed-loop run. \\
\artifact{scripts/aggregate_final_campaign.py} & Aggregates run-level KPI metrics, readiness status, paper tables, and comparison plots. \\
\artifact{data/output/final/ieee_250veh_deepc_campaign_01/01_training/deepc_dataset_txpower_count_tb0p1_dt0p5_n10/} & Final Tx-power-only \deePC{} dataset used by the controller. \\
\artifact{data/output/final/ieee_250veh_deepc_campaign_01/02_runs/fixed_10hz/run_101/} & Preliminary fixed 10 Hz baseline run used for the current comparison. \\
\artifact{data/output/final/ieee_250veh_deepc_campaign_01/02_runs/deepc_count_tb0p1_prr_only/run_931/} & Preliminary \deePC{} closed-loop run used for feasibility evidence. \\
\bottomrule
\end{longtable}

\section{Results and Evidence}

\subsection{Collected Input-Output Simulation Evidence}

The first result layer is the collected open-loop ns-3 data used to prepare
the \deePC{} dataset. In this stage, \deePC{} is not controlling the simulator;
the purpose is to collect measured input-output behavior under varying
transmit-power and traffic-density conditions. The important evidence is that
the simulation produces nonconstant PRR, PIR, and CBR values, and that these
outputs vary with both the communication input and the observed traffic
context.

\begin{figure}[htbp]
    \centering
    \includegraphics[width=0.92\linewidth]{hankel_data_quality_corrected/metric_distributions.png}
    \caption{Distribution of collected PRR, PIR, and CBR samples from the open-loop ns-3 training trace.}
    \label{fig:metric-distributions}
\end{figure}

Figure~\ref{fig:metric-distributions} shows the measured output ranges used
for data-driven modeling. This matters because a useful \deePC{} dataset should
not contain only a single operating point; it should contain enough variation
in the measured outputs to support later closed-loop optimization.

\begin{figure}[htbp]
    \centering
    \includegraphics[width=0.92\linewidth]{hankel_data_quality_corrected/tx_power_io_relationships.png}
    \caption{Collected open-loop input-output relationships between transmit power and PRR, PIR, and CBR. Red markers show mean output values at each transmit-power level.}
    \label{fig:txpower-io}
\end{figure}

Figure~\ref{fig:txpower-io} summarizes the relationship between the controlled
input, transmit power, and the measured communication outputs. The scatter is
expected because PRR, PIR, and CBR are also affected by vehicle density,
resource contention, and mobility. This is why the controller uses recent KPI
history and traffic context rather than transmit power alone.

\begin{table}[htbp]
\centering
\caption{Open-loop KPI behavior by observed traffic-density bin.}
\label{tab:density-bin-open-loop}
\begin{tabular}{lrrrr}
\toprule
Density bin & Mean active vehicles & Mean PRR & Mean PIR (s) & Mean CBR \\
\midrule
Low & 78.50 & 0.6814 & 0.2573 & 0.3497 \\
Medium & 94.32 & 0.6658 & 0.2708 & 0.3632 \\
High & 110.11 & 0.6469 & 0.2805 & 0.3753 \\
\bottomrule
\end{tabular}
\end{table}

\noindent\footnotesize Source:
\artifact{data/output/final/ieee_250veh_deepc_campaign_01/01_training/prbs_training_250veh_run001/analysis_plots/density_bin_summary.csv}.
\normalsize

Table~\ref{tab:density-bin-open-loop} shows the main density trend in the
collected data. As the number of active vehicles increases, PRR decreases while
PIR and CBR increase. This confirms that the control problem is not only a
single-input problem; network density is an important context variable for
interpreting the measured communication outputs.

\subsection{Offline Data-Window Check Evidence}

The offline data-window check shows that the collected ns-3 KPI trace contains
some input-output structure after the open-loop data-collection run. On the
current final-campaign training dataset, the linear ridge window diagnostic
achieves overall RMSE values of 0.049929 for PRR, 0.039214 for PIR, and
0.019695 for CBR.

These values support the first research milestone: the open-loop ns-3 trace can
be converted into a structured Hankel dataset for later closed-loop \deePC{}
control. However, the mixed $R^2$ values mean this offline window check
should be treated only as a data-representation diagnostic, not as evidence
that \deePC{} was controlling the open-loop simulation.

\subsection{Preliminary Closed-Loop Feasibility}

The closed-loop simulation result presented in this report is restricted to
10 Hz operation. The current comparison is preliminary because the repeated
matched campaign is not finished. The available aggregate evidence is shown in
Table~\ref{tab:closed-loop}.

\begin{table}[htbp]
\centering
\caption{Preliminary closed-loop comparison. Use for orientation only until matched repeated runs are complete.}
\label{tab:closed-loop}
\resizebox{\linewidth}{!}{%
\begin{tabular}{lrrrrrr}
\toprule
Method & Mean PRR & Mean PIR (s) & Mean CBR & CBR $>0.6$ & Success & Mean solve (s) \\
\midrule
Fixed 10 Hz, 20 dBm & 0.5748 & 0.1601 & 0.4103 & 0.0000 & -- & -- \\
\deePC{} Tx-power control & 0.5739 & 0.1603 & 0.4104 & 0.0000 & 1.0000 & 0.0199 \\
\bottomrule
\end{tabular}}
\end{table}

\noindent\footnotesize Source:
\artifact{data/output/final/ieee_250veh_deepc_campaign_01/03_analysis/tables/paper_main_table.csv}.
\normalsize

Both rows in Table~\ref{tab:closed-loop} are preliminary partial-run evidence,
not completed matched repeated-run results.

The direct comparison is \deePC{} versus fixed 10 Hz because both use 10 Hz
CAM generation. This keeps the simulation-result discussion focused on the
same nominal message-generation rate.

The preliminary \deePC{} result is very close to fixed 10 Hz in PRR, PIR, and
CBR. This means the current evidence should not be described as a demonstrated
PRR gain. The positive result is that \deePC{} actively changes transmit power
while maintaining comparable KPI behavior and solving fast enough for the
online control interval. The available run reports 285 control updates, 100\%
response success, a mean solve time of 0.0199 s, and a maximum solve time of
0.167194 s. Both timing values are below the 0.5 s control interval.

\begin{figure}[htbp]
    \centering
    \includegraphics[width=0.92\linewidth]{simulation_results_10hz/closed_loop_10hz_comparison.png}
    \caption{Preliminary 10 Hz closed-loop comparison between the fixed 10 Hz baseline and the \deePC{} transmit-power controller.}
    \label{fig:closed-loop-10hz-comparison}
\end{figure}

\begin{figure}[htbp]
    \centering
    \includegraphics[width=0.92\linewidth]{simulation_results_10hz/closed_loop_10hz_kpi_timeseries.png}
    \caption{Preliminary 10 Hz KPI time-series comparison over the overlapping fixed 10 Hz and \deePC{} run interval.}
    \label{fig:closed-loop-10hz-timeseries}
\end{figure}

Figure~\ref{fig:closed-loop-10hz-timeseries} shows the KPI behavior over time
for the overlapping interval of the available fixed 10 Hz and \deePC{} runs.
The curves support the table-level interpretation: the preliminary \deePC{}
run is close to the fixed 10 Hz baseline, so the current result is best framed
as feasibility and comparable behavior rather than proven PRR improvement.

\begin{figure}[htbp]
    \centering
    \includegraphics[width=0.92\linewidth]{simulation_results_10hz/deepc_controls.png}
    \caption{Applied transmit-power trajectory from the preliminary \deePC{} closed-loop run.}
    \label{fig:deepc-controls}
\end{figure}

Figure~\ref{fig:deepc-controls} shows that the \deePC{} controller is not only
being evaluated offline; it is actively sending transmit-power commands during
the ns-3 run. The command remains within the configured 10--23 dBm bounds.

\begin{figure}[htbp]
    \centering
    \includegraphics[width=0.92\linewidth]{simulation_results_10hz/deepc_controller_timing.png}
    \caption{Controller solve-time behavior in the preliminary \deePC{} closed-loop run.}
    \label{fig:controller-timing}
\end{figure}

Figure~\ref{fig:controller-timing} supports the real-time feasibility claim:
the mean solve time is about 0.0199 s and the maximum observed solve time is
0.167194 s, both below the 0.5 s control interval.

\noindent\footnotesize Source:
\artifact{data/output/final/ieee_250veh_deepc_campaign_01/02_runs/deepc_count_tb0p1_prr_only/run_931/bridge/controller_solve_times.csv}
and
\artifact{data/output/final/ieee_250veh_deepc_campaign_01/02_runs/deepc_count_tb0p1_prr_only/run_931/plots/summary_metrics.json}.
\normalsize

\subsection{Evidence Boundary}

The paper-readiness report explicitly marks the repeated campaign as
incomplete. For the intended matched run IDs 101--105, fixed 10 Hz currently
has only a partial run 101 with KPI data ending at 161.0 s instead of the
290 s evaluation end. The main \deePC{} method has missing run IDs 101--105,
and the current \deePC{} result comes from an extra non-matched run, run 931.

Therefore, the correct interpretation is:

\begin{quote}
The collected results validate the feasibility of \deePC-based online
transmit-power adaptation for NR-V2X and motivate matched repeated-run
evaluation.
\end{quote}

The report should not claim:

\begin{quote}
\deePC{} improves PRR over fixed 10 Hz.
\end{quote}

\section{Defensibility and Limitations}

The work is defensible as a research progress milestone because it has moved
from raw ns-3 traces to an operational closed-loop control pipeline. The
implemented pieces include data extraction, \deePC{} dataset construction,
offline data-window validation, bridge-based ns-3 integration, online optimization,
diagnostic logging, and campaign aggregation. These are substantial engineering
and research contributions even before the final statistical campaign is
complete.

The current defensible claims are:
\begin{enumerate}
    \item \deePC-ready NR-V2X datasets have been built from ns-3 NGSIM mobility traces.
    \item A linear held-out data-window diagnostic has been evaluated on PRR, PIR, and CBR using the current final-campaign training dataset.
    \item The ns-3 file bridge can apply external controller commands during a run.
    \item The Python OSQP \deePC{} controller can solve online transmit-power commands with solve times well below the 0.5 s control interval in the preliminary run.
    \item The current closed-loop evidence is comparable to fixed 10 Hz in the partial run but does not yet prove improvement.
\end{enumerate}

The main limitations are:
\begin{enumerate}
    \item The matched repeated campaign is incomplete, so confidence intervals across repeated runs are not meaningful yet.
    \item The work currently uses one processed NGSIM US-101 active-window mobility scenario.
    \item The validation is simulation-based and does not include field experiments.
    \item The current controller adapts transmit power only; it does not control SPS parameters, MCS, or resource-selection settings.
    \item Threshold-DCC should remain excluded from the main result table until the placeholder controller is replaced by a validated implementation.
\end{enumerate}

\section{Next Steps}

The most important next step is to complete the matched repeated ns-3 campaign.
The minimum professor- and paper-ready campaign should include the same run IDs
for fixed 10 Hz and the selected \deePC{} method:
\artifact{fixed_10hz} and \artifact{deepc_count_tb0p1_prr_only} for run IDs
101--105. Each run should reach the evaluation end time, normally 290 s for
the 300 s simulation with 10 s cooldown.

After completing the runs, I should regenerate the campaign analysis with:

\begin{verbatim}
./v2x_env/bin/python scripts/aggregate_final_campaign.py
\end{verbatim}

The regenerated outputs should include:
\begin{itemize}
    \item \artifact{data/output/final/ieee_250veh_deepc_campaign_01/03_analysis/paper_readiness_report.md};
    \item \artifact{data/output/final/ieee_250veh_deepc_campaign_01/03_analysis/tables/paper_main_table.csv};
    \item \artifact{data/output/final/ieee_250veh_deepc_campaign_01/03_analysis/tables/aggregate_metrics.csv};
    \item comparison plots under \artifact{data/output/final/ieee_250veh_deepc_campaign_01/03_analysis/plots/}.
\end{itemize}

Once the matched campaign is complete, the results should be checked against
the defensibility checklist before any paper-style claim is made. The main
table should report mean PRR, mean PIR, mean CBR, CBR $> 0.6$ rate, controller
success rate, and mean solve time with 95\% confidence intervals across runs.
The analysis should also include density-bin metrics and prediction-error
diagnostics, because \deePC{} may be most useful in particular traffic-density
regimes rather than uniformly across the whole scenario.

If the completed matched runs show a statistically meaningful improvement or
stability advantage, this progress report can be converted into a conference
paper. If the completed runs show similar performance to fixed 10 Hz, the paper
can still be framed around feasibility, reproducibility, and lessons learned
for data-driven NR-V2X control.

\section{Evidence and Reproducibility Pointers}

The numerical claims in this report trace to the artifacts in
Table~\ref{tab:evidence}.

\begin{longtable}{p{0.31\linewidth}p{0.62\linewidth}}
\caption{Evidence and artifact pointers.}
\label{tab:evidence}\\
\toprule
Evidence & Artifact \\
\midrule
\endfirsthead
\toprule
Evidence & Artifact \\
\midrule
\endhead
Research history and milestones & \artifact{docs/deepc_research_progress.md} \\
Conservative paper narrative & \artifact{docs/deepc_nr_v2x_conference_paper_draft.md} \\
Claim safety checklist & \artifact{docs/deepc_nr_v2x_defensibility_checklist.md} \\
Offline data-window diagnostic metrics & \artifact{data/output/final/ieee_250veh_deepc_campaign_01/01_training/open_loop_predictions/metrics.json} \\
Final campaign manifest & \artifact{data/output/final/ieee_250veh_deepc_campaign_01/campaign_manifest.json} \\
Tx-power \deePC{} dataset summary & \artifact{data/output/final/ieee_250veh_deepc_campaign_01/01_training/deepc_dataset_txpower_count_tb0p1_dt0p5_n10/dataset_summary.json} \\
Preliminary closed-loop table & \artifact{data/output/final/ieee_250veh_deepc_campaign_01/03_analysis/tables/paper_main_table.csv} \\
Run-level preliminary metrics & \artifact{data/output/final/ieee_250veh_deepc_campaign_01/03_analysis/tables/run_metrics.csv} \\
Paper readiness status & \artifact{data/output/final/ieee_250veh_deepc_campaign_01/03_analysis/paper_readiness_report.md} \\
Acceptance checklist & \artifact{data/output/final/ieee_250veh_deepc_campaign_01/04_tests/acceptance_checklist.md} \\
\bottomrule
\end{longtable}

\section{Evaluation Checklist}

\begin{itemize}
    \item ns-3 NR-V2X open-loop KPI data collected and checked.
    \item \deePC-ready dataset generated with train, validation, and test splits.
    \item Hankel matrices exported with documented dimensions.
    \item \deePC{} mathematical formulation implemented for transmit-power control.
    \item ns-3 JSON file bridge implemented for closed-loop control.
    \item Python OSQP controller implemented and connected to ns-3.
    \item Preliminary 10 Hz closed-loop result collected.
    \item Controller success rate reached 100\% in the preliminary run.
    \item Mean controller solve time was below the 0.5 s control interval.
    \item Matched repeated runs 101--105 are still incomplete.
\end{itemize}

\section{Conclusion}

This project has reached a meaningful implementation and feasibility milestone.
I have built the \deePC{} data pipeline, checked the offline data-window representation,
prepared the ns-3 closed-loop bridge, implemented online transmit-power
optimization, and generated campaign-level analysis tooling. The preliminary
closed-loop run shows that the controller can operate within the required
control interval and produce KPI behavior comparable to fixed 10 Hz operation.

The correct current conclusion is conservative: the work demonstrates a
working end-to-end \deePC{} control pipeline and preliminary real-time
feasibility for NR-V2X transmit-power adaptation. The next phase is to complete
matched repeated runs so the final paper can determine whether this feasibility
result also becomes a statistically defensible reliability or congestion-control
improvement.

\end{document}
